\section{Networks and Interconnectedness in Financial Markets}
\label{sec:ownership_networks}

The principal measures reviewed so far, from institutional ownership and crowding statistics to the fragility index of \textcite{greenwood2011stock}, share a common feature: each summarizes the ownership base of a stock through one or a small number of scalar quantities. A separate strand of the literature suggests that this compression may discard economically relevant information, since the manner in which market participants are connected can itself carry information about financial risk and its transmission.

\textcite{billio2012econometric} propose measures of financial connectedness based on principal-component analysis and Granger-causality networks constructed from the monthly returns of banks, broker-dealers, insurers, and hedge funds. They find that these sectors became highly interconnected over the preceding decade and that their connectedness measures can identify and quantify periods of financial crisis. Although this evidence concerns relationships among financial institutions rather than institutional ownership of individual stocks, it establishes a broader principle that motivates the present thesis: the topology of a financial network can contain information about systemic risk and the transmission of shocks.

Closer to the setting of this thesis, \textcite{anton2014connected} study stocks connected through common active mutual fund ownership. They show that greater shared ownership predicts higher return correlation after controlling for systematic risk exposures, style and sector similarity, and other stock-pair characteristics. Using the 2003 mutual fund trading scandal as a natural experiment, they also provide evidence that shared ownership causes this excess comovement rather than merely reflecting omitted similarities between connected firms. Their findings demonstrate that the identities of a stock's institutional owners, and particularly the connections created by common ownership, can have measurable asset-pricing consequences beyond those captured by conventional firm and stock-pair characteristics.

These findings motivate representing institutional holdings explicitly as a network rather than solely as a collection of independent stock-level indicators. In such a representation, institutional managers and securities form two distinct sets of nodes, with reported ownership positions defining the edges between them. This bipartite structure preserves the identity of each stock's owners and makes it possible to characterize how managers are connected through similar portfolios and how stocks are exposed to those connections. The resulting structure can be summarized through engineered features describing owner similarity, owner centrality, and concentration across manager communities, or it can be supplied directly to models that learn from the graph itself.